Table~\ref{tab:related} summarises explicitly reported features of selected related studies and the present study. These feature indicators describe study scope rather than comparative performance.

\begingroup
\footnotesize
\setlength{\tabcolsep}{2pt}
\renewcommand{\arraystretch}{1.4}
\begin{longtable}{@{}>{\raggedright\arraybackslash}p{0.18\linewidth}*{7}{>{\centering\arraybackslash}p{0.105\linewidth}}@{}}
\caption{Reported features of selected liver transcriptomic studies and the present study}\label{tab:related}\\
\toprule
Study & Multiple cohorts & Separate F0--F4 analysis & Gene stage-profile clustering & Adjacent-stage contrasts & Curated knowledge graph & Multi-database drug data & Independent-cohort or experimental follow-up \\
\midrule
\endfirsthead
\toprule
Study & Multiple cohorts & Separate F0--F4 analysis & Gene stage-profile clustering & Adjacent-stage contrasts & Curated knowledge graph & Multi-database drug data & Independent-cohort or experimental follow-up \\
\midrule
\endhead
\bottomrule
\endfoot
Hoang et al., 2019 \cite{ref4} & $\times$ & $\checkmark$ & $\times$ & $\times$ & $\times$ & $\times$ & $\times$ \\
Govaere et al., 2020 \cite{govaere2020} & $\checkmark$ & $\checkmark$ & $\times$ & $\times$ & $\times$ & $\times$ & $\checkmark$ \\
Piras and DiStefano, 2024 \cite{piras2024} & $\checkmark$ & $\times$ & $\times$ & $\times$ & $\times$ & $\times$ & $\checkmark$ \\
Yang et al., 2026 \cite{yang2026} & $\checkmark$ & $\times$ & $\times$ & $\times$ & $\times$ & $\times$ & $\checkmark$ \\
Maiers et al., 2026 \cite{maiers2026} & $\checkmark$ & $\times$ & $\times$ & $\times$ & $\times$ & $\times$ & $\checkmark$ \\
\midrule
\textbf{Present study} & $\checkmark$ & $\checkmark$ & $\checkmark$ & $\checkmark$ & $\checkmark$ & $\checkmark$ & $\times$ \\
\end{longtable}
\textbf{Reading the symbols.} $\checkmark$: explicitly documented in the reviewed study descriptions. $\times$: not documented there for the precisely defined feature; this is not proof that an unreported analysis was never performed. Multiple cohorts includes discovery/replication designs and does not imply statistical meta-analysis. Gene stage-profile clustering specifically means clustering genes using their F0--F4 expression profiles; gene shaving, sample clustering and coexpression-module construction are different procedures. Piras and DiStefano principally compared diagnostic groups; Maiers et al. compared F3--F4 with F0--F1, excluding F2. Protein-interaction and regulatory networks are not counted as curated knowledge-graph integration. The present resource integrates OptimusKG with drug data from ChEMBL, DGIdb, GtoPdb and Open Targets; this feature does not imply validated drug prediction.

\textbf{Robustness and validation.} Hoang et al. used cross-validation, which is internal assessment. Govaere et al. reported independent-cohort replication, serum-protein measurements and selected functional experiments. Piras and DiStefano evaluated module preservation and module differential expression in cohorts separate from their module-discovery cohort; this does not constitute independent validation of their entire meta-analysis or causal network. Yang et al. reported RT-qPCR and selected protein measurements. Maiers et al. confirmed seven of their top 50 selected genes in an independent cohort. The present study used leave-one-study-out refits, clustering resampling and sensitivity analyses; independent-cohort and experimental validation remain absent.
\endgroup

